\documentclass[11pt]{article}
\usepackage[T1]{fontenc}
\usepackage{lmodern}
\usepackage[margin=1in]{geometry}
\usepackage{amsmath,amssymb}
\usepackage{microtype}
\usepackage{booktabs}
\usepackage{xcolor}
\usepackage[colorlinks=true,linkcolor=blue!50!black]{hyperref}

\newcommand{\key}[1]{\par\smallskip\noindent{\bfseries #1}\par\nobreak\vspace{-2pt}}
\setlength{\abovedisplayskip}{4pt}\setlength{\belowdisplayskip}{4pt}
\setlength{\abovedisplayshortskip}{2pt}\setlength{\belowdisplayshortskip}{2pt}

\title{\bfseries Step 5 on one example\\[4pt]
\large\normalfont reading one documented pass back: the three states, the ledger, and one margin}
\author{}\date{}

\begin{document}
\maketitle
\vspace{-2.6em}

\noindent
Steps 1--4 are implications: a region maps to a smaller region, or to no solution. Step 5 is the
statement about a \emph{run} --- that a documented execution visited every region and left none
undecided --- and it is the step that makes the theorem conditional. Its tests are therefore not
inequalities on matrices but three questions about a ledger: does every record end in exactly one of
three states, is the list of unresolved records empty, and is every threshold comparison made with a
margin that survives binary64?

\smallskip
\noindent
The example input is a real one: the archived record census of a Stage-C pass of the documented run,
$23{,}506$ graph records produced by $92$ commands. It is small enough to recount and rich enough to
show all three states --- including the one that is not empty.

\key{1.\ The three states, counted}
Every object the earlier stages produce carries a record, and a record ends in exactly one of
\emph{excluded} (a necessary condition failed, with the margin of Appendix~A), \emph{resolved} (its
completions were transferred to a certified ball, later excluded), or \emph{unresolved} (neither).
In this pass the census is
\begin{center}\small
\begin{tabular}{@{}lrrrrr@{}}
\toprule
state & records & min steps & median steps & max steps & total steps\\
\midrule
excluded \ (\texttt{nopattern}: completed negative search) & $23{,}343$ & $76$ & $38{,}453$ & $4{,}986{,}110$ & $2{,}932{,}018{,}956$\\
resolved \ (\texttt{positive}: a pattern found) & $0$ & --- & --- & --- & ---\\
unresolved \ (\texttt{cap}: stopped at $5{,}000{,}001$) & $163$ & $5{,}000{,}001$ & $5{,}000{,}001$ & $5{,}000{,}001$ & $815{,}000{,}163$\\
\midrule
all records & $\mathbf{23{,}506}$ & $76$ & --- & $5{,}000{,}001$ & $\mathbf{3{,}747{,}019{,}119}$\\
\bottomrule
\end{tabular}
\end{center}
Three things in that table are worth reading twice. Every completed search ended \emph{negative} ---
$23{,}343$ exclusions and no pattern --- which is what a theorem of this shape predicts. The capped
records stopped at exactly $5{,}000{,}001$ steps, one past the cap, so the cap is what stopped them and
not the problem. And the largest completed search took $4{,}986{,}110$ steps, $99.7\%$ of the cap:
the cap is not far above the hardest genuine search in this pass, which is the quantitative version of
``the fallback matters''.

\key{2.\ The ledger, and the one line that fails}
The identity a completed execution must satisfy is
\begin{equation}\label{eq:ledger}
\text{records}=\text{excluded}+\text{resolved}+\text{unresolved},\qquad\text{unresolved}=0 .
\end{equation}
Recounted from the records rather than from the log:
\[
23{,}506=23{,}343+0+163 \quad\checkmark
\qquad\text{but}\qquad
\text{unresolved}=163\neq0 .
\]
The first half is the accounting: the three states partition the records, with nothing double-counted
and nothing lost. The second half is the substantive condition, and it \emph{fails}. This pass is
therefore not a completed execution, and no statement about four bases follows from it: $163$ regions
are simply undecided and would have to be refined and reprocessed. That is exactly the right outcome
for a pass whose cap was $5\cdot10^{6}$ steps, and saying so is the whole discipline of the step: a cap
that runs out is an admission, never an exclusion.

\key{3.\ Three records, as miniatures}
A table of the extremes is enough to see what a record is:
\begin{center}\small
\begin{tabular}{@{}llrrrrr@{}}
\toprule
record & state & bins & balls & $O$ edges & $U$ edges & steps\\
\midrule
\texttt{final0/230} \ (fewest steps) & excluded & $1{,}535$ & $52$ & $122$ & $384$ & $76$\\
\texttt{batch9/140} \ (median) & excluded & $25{,}034$ & $268$ & $2{,}981$ & $15{,}981$ & $38{,}852$\\
\texttt{batch10/43} \ (capped) & unresolved & $31{,}852$ & $318$ & $10{,}399$ & $19{,}178$ & $5{,}000{,}001$\\
\bottomrule
\end{tabular}
\end{center}
The three columns before the step count are the sizes of the objects Step~4 built --- the final bin
layer, the vector balls, and the two graphs --- so a record is a complete, checkable summary of one
region: how it was cut, what was kept, and what the search concluded. Across the pass there were
$5{,}272{,}329$ vector balls and $368{,}119{,}692$ edges.

\key{4.\ One margin, computed}
Each node of the decision tree carries an inequality, and each inequality is evaluated in binary64. The
contract is that every threshold is compared with a margin larger than the rounding. Take the smallest
constant in the proof, the vanishing-sum slack of Step~1, $5\pi w^2/4$ at $w=1/112$:
\[
\frac{5\pi}{4\cdot112^2}=\frac{5\pi}{50176}=3.13057303650130863\times10^{-4},
\qquad
\text{fl}\Bigl(\frac{5\pi}{50176}\Bigr)=0.00031305730365013086 .
\]
The binary64 value is wrong by $0$ ulps here; a five-rounding bound $k\,2^{-53}$ gives a relative
width $5.6\times10^{-16}$ and an absolute width $1.7\times10^{-19}$. The code, by contrast, compares
its thresholds with an absolute slack of $2\times10^{-7}$ --- larger than the outward interval by a
factor $1.15\times10^{12}$. The same computation for the support bound $1.25\pi/N$ with $N=112$:
$0.0350624180088$, outward width $1.9\times10^{-17}$.

So on this example floating point is not where the risk is: the margins the implementation uses are
twelve orders of magnitude larger than the arithmetic error they have to cover. The risk that remains
is the other kind --- whether the \emph{constant} in a test is the right one --- and that is a
mathematical question about the paper, not a numerical one about the machine.

\key{5.\ The trap in the run's own log}
The pass's log line reads
\[
\texttt{done}=23{,}506,\quad \texttt{bad}=163,\quad \texttt{graph}=0,\quad
\texttt{cap}=163,\quad \texttt{over}=0,\quad \texttt{cluster}=0 .
\]
Read quickly, \texttt{graph}=0 says ``no graph search ran''. It is false: $23{,}506$ graph searches ran,
with $3{,}747{,}019{,}119$ steps between them. The counter prints the number of searches that
\emph{found a candidate pattern} --- the implementation's \texttt{positive} --- so $0$ means ``none
found one'', which is the same statement as $\text{resolved}=0$ in \S1. It is the single easiest way to
misread this run, and it is why the census above was recomputed from the records rather than taken from
the log.

\key{6.\ What this example does not say}
It is one pass of one stage of the pipeline, and the pipeline has four stages; a completed execution
needs every stage and every region, and the other passes are not recounted here. It does not verify the
implementation: the numbers are read from records the programs wrote, and a record that was written
wrongly would be counted wrongly. It says nothing about compiler or hardware behaviour beyond the one
archived run. And the ledger checks that the records partition, not that each exclusion is correct ---
that is the job of Steps 1--4, one inequality at a time.

\medskip
\noindent\footnotesize\textit{Provenance.} The three states, the condition \eqref{eq:ledger} and the
arithmetic contract are Step~5's, in \texttt{mub-step5.pdf} and preprint \S7 and Appendix~A. The
records are \texttt{certdata/runs/h0-stagec-probe.json}, the archive of the pass; the log line is
\texttt{h0-stagec-probe-summary.json}. Every number above was recomputed by
\texttt{examples/step5\_example.py}, which reads the archive and prints the table.

\end{document}
